\chapter{Nền tảng hệ hỗ trợ quyết định}

\section{Từ dữ liệu đến hành động}

Dữ liệu là các quan sát thô, chẳng hạn nội dung một tweet, một ảnh, thời điểm
đăng và tên sự kiện. Khi dữ liệu được làm sạch, đặt trong bối cảnh và tổng hợp,
nó trở thành \emph{thông tin}: bài đăng có khả năng hữu ích, liên quan đến hạ
tầng và có mức bất đồng text--ảnh cao. Khi thông tin được kết hợp với quy tắc
nghiệp vụ, giới hạn nguồn lực và trách nhiệm của người sử dụng, hệ thống mới
có thể tạo \emph{khuyến nghị hành động}: chuyển bài đăng tới đội hạ tầng,
xếp Medium và yêu cầu giám sát viên kiểm tra.

\begin{figure}[H]
\centering
\begin{tikzpicture}[
  node distance=7mm,
  stage/.style={
    draw=HUSTBlue,
    rounded corners=2pt,
    minimum width=3.1cm,
    minimum height=1.05cm,
    align=center,
    fill=boxgray
  },
  arrow/.style={-{Latex[length=2.5mm]},thick,HUSTBlue}
]
\node[stage] (data) {Dữ liệu\\tweet, ảnh, nhãn};
\node[stage,right=of data] (info) {Thông tin\\xác suất, category};
\node[stage,right=of info] (knowledge) {Tri thức/chính sách\\severity, capacity};
\node[stage,right=of knowledge] (action) {Khuyến nghị\\priority, routing};
\draw[arrow] (data) -- (info);
\draw[arrow] (info) -- (knowledge);
\draw[arrow] (knowledge) -- (action);
\end{tikzpicture}
\caption{Chuỗi chuyển đổi từ dữ liệu thô tới khuyến nghị hành động.}
\end{figure}

Việc phân biệt bốn mức trên giúp tránh một sai lầm phổ biến: một xác suất dự
báo cao không tự nó là quyết định. Ví dụ, xác suất 0,82 cho lớp hạ tầng chưa
cho biết đội nào tiếp nhận hoặc trường hợp đó có khẩn cấp hơn một lời kêu cứu
khác hay không. Những câu hỏi đó cần thêm chính sách và bối cảnh vận hành.

\section{Khái niệm hệ hỗ trợ quyết định}

Hệ hỗ trợ quyết định là hệ thống thông tin kết hợp dữ liệu, mô hình phân tích
và giao diện tương tác để giúp con người giải quyết các vấn đề mà việc lựa
chọn không thể được tự động hóa hoàn toàn \cite{power2002,slides}. DSS không
nhất thiết phải đưa ra một đáp án duy nhất. Giá trị của nó có thể nằm ở việc:

\begin{itemize}[leftmargin=*]
  \item giảm khối lượng thông tin con người phải đọc;
  \item cung cấp nhiều góc nhìn và chỉ số có thể kiểm tra;
  \item sắp xếp các phương án theo tiêu chí rõ ràng;
  \item cảnh báo trường hợp bất thường hoặc thiếu chắc chắn;
  \item ghi lại lý do của khuyến nghị để người sử dụng chịu trách nhiệm cuối.
\end{itemize}

Trong dự án này, DSS không thay thế trung tâm điều phối. Nó biến hàng nghìn
bài đăng thành một hàng đợi có cấu trúc, nhưng con người vẫn kiểm tra vị trí,
độ xác thực và điều kiện hiện trường trước khi triển khai nguồn lực.

\section{Mức độ cấu trúc của quyết định}

Một \textbf{quyết định có cấu trúc} có quy tắc và đầu vào rõ ràng, có thể thực
hiện lặp lại. Ví dụ, nếu một bài đăng đã được xác nhận thuộc nhóm hư hỏng hạ
tầng thì chuyển về Infrastructure Team là một quy tắc tương đối có cấu trúc.

Một \textbf{quyết định bán cấu trúc} có một phần được mô hình hoặc quy tắc hỗ
trợ nhưng vẫn cần phán đoán con người. Xếp mức ưu tiên cho bài đăng là quyết
định bán cấu trúc: Risk Score tạo thứ tự ban đầu, nhưng điều phối viên có thể
thay đổi khi biết bệnh viện nào đang quá tải hoặc tuyến đường nào đã được xử
lý.

Một \textbf{quyết định phi cấu trúc} phụ thuộc mạnh vào kinh nghiệm và bối
cảnh khó biểu diễn thành quy tắc cố định, chẳng hạn quyết định điều động lực
lượng ngoài đời trong một thảm họa đang diễn biến. Phạm vi dự án dừng ở mức
hỗ trợ cho loại quyết định này, không tự động thực thi.

\section{Chu trình ra quyết định}

Theo mô hình của Simon, quá trình ra quyết định có thể được mô tả qua các pha
nhận biết vấn đề, thiết kế phương án, lựa chọn và thực thi
\cite{simon1960}. Áp dụng vào dự án:

\begin{enumerate}[leftmargin=*]
  \item \textbf{Nhận biết (Intelligence):} thu nhận bài đăng, loại dữ liệu lỗi,
        phát hiện bài informative và nhận diện loại nhu cầu.
  \item \textbf{Thiết kế (Design):} tạo các phương án xử lý như Emergency,
        Relief, Infrastructure, Coordination hoặc Manual Review.
  \item \textbf{Lựa chọn (Choice):} dùng Risk Score, Priority và routing để đề
        xuất phương án phù hợp.
  \item \textbf{Thực thi và phản hồi:} người điều phối xác nhận, thực hiện và
        ghi nhận kết quả để sau này hiệu chỉnh chính sách.
\end{enumerate}

Dữ liệu CrisisMMD chỉ cung cấp ground truth cho informative và humanitarian,
không có phản hồi thực tế về Priority hoặc đội được điều động. Vì vậy dự án
thực hiện tốt ba pha đầu ở mức nguyên mẫu; pha phản hồi vận hành được xác định
là hướng phát triển.

\section{Ba nhóm năng lực của hệ thống}

\subsection{Phân tích mô tả và chẩn đoán}

Phân tích mô tả trả lời ``dữ liệu đang có đặc điểm gì?'': bao nhiêu bài theo
sự kiện, tỷ lệ informative, phân bố tám lớp và độ dài tweet. Phân tích chẩn
đoán đi thêm một bước để tìm lý do hoặc mối quan hệ: lớp hiếm, sự kiện có phân
bố khác nhau, text và ảnh bất đồng ở nhóm nào. EDA, K-Means và Apriori thuộc
nhóm này.

\subsection{Phân tích dự báo}

Phân tích dự báo trả lời ``mẫu mới có khả năng thuộc lớp nào?''. Các classifier
văn bản và ảnh tạo ra xác suất informative và xác suất của tám nhóm
humanitarian. Kết quả này được đánh giá bằng nhãn chính thức, confusion matrix
và các chỉ số phù hợp với dữ liệu mất cân bằng.

\subsection{Phân tích định hướng hành động}

Phân tích định hướng hành động trả lời ``nên làm gì với dự báo?''. Tầng DSS
dùng xác suất, mức nghiêm trọng, từ khóa và conflict score để đề xuất thứ tự
ưu tiên, đội tiếp nhận và nhu cầu kiểm duyệt. Vì CrisisMMD không có nhãn
Priority, đây là chính sách minh bạch chứ không phải một mô hình được chứng
minh tối ưu.

\begin{table}[H]
\centering
\small
\caption{Phân biệt các nhóm năng lực phân tích trong dự án.}
\begin{tabularx}{\textwidth}{p{3.1cm}p{4.3cm}XX}
\toprule
\textbf{Nhóm} & \textbf{Câu hỏi} & \textbf{Đầu ra} & \textbf{Thành phần}\\
\midrule
Mô tả/chẩn đoán & Dữ liệu có đặc điểm và vấn đề gì? &
Phân bố, insight, cụm, luật kết hợp & EDA, K-Means, Apriori\\
Dự báo & Bài đăng có khả năng thuộc lớp nào? &
Xác suất informative và humanitarian & Sáu classifier, Late Fusion\\
Định hướng hành động & Nên xử lý bài đăng như thế nào? &
Risk, Priority, Routing, Review & Rule engine và dashboard\\
\bottomrule
\end{tabularx}
\end{table}

\section{Kiến trúc DSS của dự án}

\begin{figure}[H]
\centering
\begin{tikzpicture}[
  node distance=7mm and 9mm,
  layer/.style={
    draw=HUSTBlue,
    rounded corners=2pt,
    text width=3.2cm,
    minimum height=1.3cm,
    align=center,
    fill=boxgray
  },
  arrow/.style={-{Latex[length=2.5mm]},thick,HUSTBlue},
  feedback/.style={-{Latex[length=2.5mm]},thick,dashed,HUSTRed}
]
\node[layer] (source) {Dữ liệu\\CrisisMMD text + image};
\node[layer,right=of source] (feature) {Biểu diễn\\TF--IDF + CLIP cố định};
\node[layer,right=of feature] (predict) {Dự báo\\text/image classifiers};
\node[layer,below=of predict] (fusion) {Hợp nhất\\Late Fusion + conflict};
\node[layer,left=of fusion] (policy) {Chính sách DSS\\Risk + Priority + Routing};
\node[layer,left=of policy] (ui) {Giao diện\\Dashboard + hàng đợi};
\draw[arrow] (source) -- (feature);
\draw[arrow] (feature) -- (predict);
\draw[arrow] (predict) -- (fusion);
\draw[arrow] (fusion) -- (policy);
\draw[arrow] (policy) -- (ui);
\draw[feedback] (ui.north) -- (source.south);
\end{tikzpicture}
\caption{Kiến trúc logic của hệ hỗ trợ quyết định đa phương thức.}
\end{figure}

Đường nét đứt biểu diễn vòng phản hồi dự kiến: quyết định của người vận hành
và kết quả xác minh có thể được lưu lại để theo dõi sai lệch, cập nhật chính
sách hoặc tái huấn luyện mô hình ở các phiên bản sau.

Kiến trúc gồm bốn khối:

\begin{enumerate}[leftmargin=*]
  \item \textbf{Dữ liệu và biểu diễn:} làm sạch tweet, tạo TF--IDF; chuẩn hóa
        ảnh và dùng CLIP tiền huấn luyện với trọng số giữ nguyên để tạo vector.
  \item \textbf{Mô hình dự báo:} các classifier cổ điển được huấn luyện trên
        CrisisMMD cho hai nhiệm vụ informative và humanitarian.
  \item \textbf{Hợp nhất và chính sách:} xác suất hai phương thức được căn
        chỉnh, kết hợp và chuyển thành khuyến nghị.
  \item \textbf{Giao diện:} người dùng xem KPI, lọc theo sự kiện, kiểm tra một
        ca và tải hàng đợi xử lý.
\end{enumerate}

\section{Vai trò của con người trong hệ thống}

Human-in-the-loop là thiết kế trong đó con người được đưa vào những điểm mà
mô hình không đủ dữ liệu, thiếu chắc chắn hoặc quyết định có hậu quả lớn. Với
dự án này, con người xuất hiện ở ba mức:

\begin{itemize}[leftmargin=*]
  \item \textbf{xác minh dữ liệu:} kiểm tra bài đăng, ảnh, vị trí và nguồn;
  \item \textbf{xác nhận khuyến nghị:} chấp nhận hoặc thay đổi Priority và
        Routing theo tình hình thực tế;
  \item \textbf{giám sát hệ thống:} theo dõi sai lệch theo sự kiện, công suất
        hàng review và hiệu quả sau triển khai.
\end{itemize}

Manual Review (kiểm tra thủ công) có một quy trình cụ thể chứ không chỉ là một
cái cờ bật/tắt. Khi hệ thống đánh dấu một tweet cần review, năm bước sau được
thực hiện bởi con người:

\begin{figure}[H]
\centering
\begin{tikzpicture}[
  font=\small,
  node distance=4mm,
  mrstep/.style={draw=HUSTBlue,rounded corners=2pt,fill=boxgray,align=center,
    text width=2.45cm,minimum height=1.35cm},
  arrow/.style={-{Latex[length=2.2mm]},thick,HUSTBlue}
]
\node[mrstep] (s1) {1. Hệ thống\\nêu \emph{lý do}\\(conflict cao)};
\node[mrstep,right=of s1] (s2) {2. Người xem\\tweet, ảnh,\\nguồn, thời điểm};
\node[mrstep,right=of s2] (s3) {3. Đối chiếu\\nguồn thứ hai\\(hotline, báo)};
\node[mrstep,right=of s3] (s4) {4. Chấp nhận\\hoặc sửa\\category \& priority};
\node[mrstep,right=of s4] (s5) {5. Ghi lại\\kết quả\\xác minh};
\draw[arrow] (s1) -- (s2);
\draw[arrow] (s2) -- (s3);
\draw[arrow] (s3) -- (s4);
\draw[arrow] (s4) -- (s5);
\end{tikzpicture}
\caption{Quy trình Manual Review: hệ thống chỉ nêu lý do, con người mới là người
xác minh và quyết định.}
\label{fig:manual-review}
\end{figure}

Manual Review vì vậy không chỉ là một ``nhãn lỗi''. Nó là một cơ chế quản trị
rủi ro:
hệ thống chủ động chuyển các trường hợp bất đồng mạnh cho người có trách nhiệm.
Tuy nhiên, công suất con người là hữu hạn. Nếu ngưỡng quá thấp, gần như mọi
bài đăng đều bị chuyển review và DSS không còn giảm tải. Vì vậy ngưỡng conflict
được chọn trên dev dưới một giới hạn công suất, sau đó mới khóa để đánh giá.

\section{Nguyên tắc phân biệt bằng chứng và chính sách}

\begin{insight}[Ba mức khẳng định trong báo cáo]
\textbf{Bằng chứng dữ liệu} là đại lượng đo được, chẳng hạn số mẫu, phân bố
nhãn, metric hoặc tỷ lệ bất đồng. \textbf{Suy luận phân tích} là cách diễn giải
có giới hạn từ bằng chứng, chẳng hạn ``mất cân bằng có thể làm Accuracy che
giấu lỗi lớp hiếm''. \textbf{Chính sách vận hành} là lựa chọn của tổ chức khi
dữ liệu không có ground truth, chẳng hạn trọng số severity hoặc ngưỡng
Priority. Ba mức này phải được trình bày tách biệt.
\end{insight}

\enlargethispage{3\baselineskip}
Nguyên tắc trên đặc biệt quan trọng với Risk Score. Dự án có thể chứng minh
classifier được đánh giá trên nhãn informative/humanitarian, nhưng không thể
chứng minh Priority là tối ưu khi chưa có nhãn Priority. Thay vào đó, báo cáo
phải công khai công thức, thử nhiều ngưỡng và mô tả điều kiện cần để xác nhận
chính sách trong tương lai.

Bảng~\ref{tab:evidence-policy} áp dụng nguyên tắc này cho một ca cụ thể, tách
rõ đâu là số đo, đâu là suy luận, đâu là lựa chọn của tổ chức và đâu là việc
nằm ngoài hệ thống.

\begin{table}[H]
\centering
\small
\caption{Bốn mức khẳng định cho một tweet ``cầu sập, người mắc kẹt''.}
\label{tab:evidence-policy}
\begin{tabularx}{\textwidth}{p{3.4cm}X}
\toprule
\textbf{Mức} & \textbf{Ví dụ trên ca này}\\
\midrule
Bằng chứng dữ liệu & Xác suất informative 0,93; nhánh ảnh và văn bản cùng cho
nhóm hạ tầng; conflict score thấp.\\
Suy luận phân tích & ``Hai phương thức đồng thuận và xác suất cao nên ca này ít
khả năng là báo động giả.''\\
Chính sách vận hành & Trọng số severity của nhóm hạ tầng và ngưỡng Priority do
nhóm chọn $\rightarrow$ xếp High.\\
Hành động ngoài hệ thống & Điều xe, gọi đội cứu hộ: \emph{không} thuộc hệ thống;
do con người và cơ quan nghiệp vụ thực hiện sau khi xác minh.\\
\bottomrule
\end{tabularx}
\end{table}

\section{Bản đồ kiến thức học phần}
\label{sec:concept-map}

Báo cáo dùng nhiều kỹ thuật, nhưng không phải kỹ thuật nào cũng là \emph{trọng
tâm} của học phần Hệ hỗ trợ quyết định. Bảng~\ref{tab:concept-map} phân biệt
phần thuộc nội dung môn học (trung tâm) với phần công cụ bổ trợ nằm ngoài slide.

\begin{table}[H]
\centering
\small
\caption{Bản đồ kiến thức: nội dung học phần và vai trò trong dự án.}
\label{tab:concept-map}
\begin{tabularx}{\textwidth}{p{2.4cm}p{5.2cm}X}
\toprule
\textbf{Nhóm} & \textbf{Nội dung học phần} & \textbf{Vai trò trong dự án}\\
\midrule
Trung tâm & DSS, quy trình ra quyết định & Từ dự báo tới hàng đợi hành động\\
Trung tâm & Học máy và đánh giá & Train/dev/test, metric, baseline\\
Trung tâm & Decision Tree, NB, k-NN, SVM, Ensemble & Sáu mô hình so sánh\\
Trung tâm & Text Mining, Social Media Mining & Làm sạch tweet, TF--IDF\\
Trung tâm & Phân cụm (Clustering) & K-Means khám phá chủ đề\\
Trung tâm & Phân tích luật kết hợp & Apriori trên từ khóa/hashtag\\
Trung tâm & BI/Dashboard & Streamlit trình bày và thao tác hàng đợi\\
\midrule
Bổ trợ & CLIP/ViT & Trích đặc trưng ảnh cố định (frozen)\\
Bổ trợ & SHA/pHash, t-SNE, bootstrap & Kiểm soát rò rỉ và độ bền\\
\bottomrule
\end{tabularx}
\end{table}

\begin{insight}[Nội dung học phần thể hiện ở đâu?]
Phần \emph{trung tâm} của bài tập là tầng DSS (Chương~\ref{chap:dss}) và sáu
classifier cùng cách đánh giá (Chương~\ref{chap:mohinh}), dựa trên Text Mining,
Clustering và Association Analysis ở Chương~\ref{chap:eda}. CLIP \emph{không}
được coi là đóng góp thuật toán chính: nó chỉ là một bộ trích đặc trưng ảnh dùng
sẵn. Điều được chấm trong học phần là cách chọn mô hình, đánh giá trung thực và
biến kết quả thành quyết định, không phải việc gọi một mô hình ảnh có sẵn.
\end{insight}

\subsection{Vai trò của các công cụ tính toán}

Một số người đọc có thể thắc mắc vì sao một dự án ``khá lớn'' lại không dùng các
framework học sâu hay xử lý phân tán. Lý do nằm ở \emph{quy mô và mục tiêu}:

\begin{itemize}[leftmargin=*]
  \item \textbf{PyTorch} chỉ đóng vai trò \emph{runtime} để nạp và chạy CLIP đã
        huấn luyện sẵn. Dự án không tự xây dựng hay fine-tune mạng nơ-ron sâu
        nào; mọi phần học máy ``của dự án'' nằm ở các classifier cổ điển trên
        scikit-learn.
  \item \textbf{PySpark} không được dùng vì 18.082 dòng và embedding 512 chiều
        vừa khít bộ nhớ một máy. Spark giải quyết bài toán dữ liệu phân tán hàng
        triệu đến hàng tỷ bản ghi; với quy mô này nó chỉ thêm chi phí vận hành mà
        không gỡ được nút thắt nào. Spark sẽ cần thiết nếu sau này dữ liệu chảy
        liên tục, đạt hàng triệu bản ghi hoặc phải xử lý phân tán nhiều nguồn.
  \item \textbf{scikit-learn, pandas} là phần lõi cho mô hình hóa và xử lý bảng;
        \textbf{Streamlit} dựng dashboard. Vai trò chi tiết của từng thư viện
        nằm ở bảng công nghệ Chương~\ref{chap:dashboard}.
\end{itemize}

Tóm lại, lựa chọn công cụ bám theo nguyên tắc tối thiểu cần thiết: chỉ thêm một
công nghệ khi nó giải quyết một nút thắt thật, không thêm để ``cho oai''.
